\subsection{Desarrollo del proyecto}

El desarrollo del proyecto se organiza mediante una secuencia
estructurada de actividades, divididas en tres fases metodológicas
principales. Cada fase está diseñada para asegurar el cumplimiento
progresivo de los objetivos específicos planteados en la
investigación, partiendo desde el diagnóstico hasta la comprobación
cuantitativa de los resultados.

\bigskip
\noindent\textbf{Fase 1: Análisis del flujo actual de información}

\smallskip
Esta fase tiene como propósito comprender la dinámica de trabajo
existente en el centro odontológico, identificar las deficiencias
operativas y establecer los lineamientos normativos obligatorios
previos a cualquier propuesta de rediseño. Dado que la medición de
tiempos de línea base forma parte integral del diagnóstico, la
toma de tiempos pre-implementación se ejecuta también dentro de
esta fase, antes de que cualquier intervención tecnológica altere
el flujo natural de trabajo. Las actividades correspondientes al
primer objetivo específico se detallan en la
Tabla~\ref{tab:actividades_fase1}.

\begin{table}[H]
\centering
\begin{tabular}{|p{4.5cm}|p{9cm}|}
\hline
\textbf{Actividad} & \textbf{Descripción} \\ \hline
Medición inicial
(pre-implemen\-tación)
  & Registro cronometrado del tiempo del ciclo administrativo
    completo mediante selección aleatoria de ciclos de atención,
    calculando el tiempo promedio y la dispersión como valores
    de línea base previos a toda intervención. \\ \hline
Revisión documental
  & Análisis de los formatos físicos y digitales vigentes
    (historias clínicas, registros de pago y recetas) para
    establecer el estado inicial de la gestión de datos. \\ \hline
Entrevistas individuales
  & Ejecución de entrevistas estructuradas con el personal
    involucrado (odontólogo general, especialista y secretaria)
    para recolectar información sobre sus rutinas y dificultades
    diarias. \\ \hline
Observación directa
  & Seguimiento de una jornada continua de atención para
    registrar el orden cronológico de las tareas, su duración
    y las interrupciones en el flujo de trabajo. \\ \hline
Revisión de normativa
  & Análisis de las regulaciones del Ministerio de Salud
    Pública referentes a la estructura y manejo obligatorio
    de las historias clínicas odontológicas, incluyendo el
    Formulario~033 y el Acuerdo Ministerial 5216-A. \\ \hline
Elaboración del mapa
de procesos
  & Construcción del esquema gráfico que representa la
    circulación actual de la información del paciente durante
    el ciclo completo de atención, identificando los puntos
    de ruptura operativa. \\ \hline
\end{tabular}
\caption{Actividades de la Fase 1: Análisis del flujo actual
         de información.}
\label{tab:actividades_fase1}
\end{table}

\noindent\textbf{Resultado de la Fase 1:} Un documento de
diagnóstico que consolida los procesos actuales, los puntos de
ruptura operativa, los requerimientos funcionales y normativos
que guiarán el diseño del sistema, y los valores de tiempo
promedio por ciclo de atención registrados como línea base.

\bigskip
\noindent\textbf{Fase 2: Diseño del modelo de información}

\smallskip
A partir de los hallazgos de la fase de diagnóstico, se
estructurará el funcionamiento del nuevo sistema y la organización
de los datos, garantizando una operación fluida e integrada. Las
actividades correspondientes al segundo objetivo específico se
detallan en la Tabla~\ref{tab:actividades_fase2}.

\begin{table}[H]
\centering
\begin{tabular}{|p{4.5cm}|p{9cm}|}
\hline
\textbf{Actividad} & \textbf{Descripción} \\ \hline
Definición de
responsabilidades
  & Determinación exacta de los datos que debe ingresar cada
    actor en las distintas etapas de atención, asegurando la
    disponibilidad de la información sin requerir traslados
    manuales entre herramientas. \\ \hline
Establecimiento de
reglas de negocio
  & Definición secuencial de los procesos (agendamiento,
    apertura de historia clínica, emisión de recetas y cobro),
    incluyendo las reglas del Formulario~033, garantizando la
    conexión automática de estos eventos dentro del sistema. \\ \hline
Descripción de módulos
funcionales
  & Especificación de las partes funcionales del sistema,
    detallando el propósito de cada módulo, la información de
    entrada requerida y los resultados generados. \\ \hline
Verificación de
cumplimiento normativo
  & Revisión del diseño para asegurar la inclusión de los
    formatos y campos exigidos por el Ministerio de Salud
    Pública, garantizando la validez legal del sistema desde
    su concepción. \\ \hline
\end{tabular}
\caption{Actividades de la Fase 2: Diseño del modelo de
         información.}
\label{tab:actividades_fase2}
\end{table}

\noindent\textbf{Resultado de la Fase 2:} Un diseño documentado
y validado del sistema de información que servirá como guía
estructural definitiva antes de iniciar la etapa de construcción
tecnológica.

\bigskip
\noindent\textbf{Fase 3: Implementación y evaluación de impacto}

\smallskip
En la fase final se construirá la herramienta tecnológica, se
desplegará en el entorno real de trabajo y se evaluará
cuantitativamente su efecto sobre la variable dependiente: el
tiempo promedio por ciclo completo de atención. Las actividades
correspondientes al tercer objetivo específico se detallan en la
Tabla~\ref{tab:actividades_fase3}.

\begin{table}[H]
\centering
\begin{tabular}{|p{4.5cm}|p{9cm}|}
\hline
\textbf{Actividad} & \textbf{Descripción} \\ \hline
Construcción e
instalación del sistema
  & Desarrollo del sistema informático en apego estricto al
    diseño validado en la Fase~2 e instalación en los equipos
    del consultorio. \\ \hline
Capacitación del
personal
  & Entrenamiento práctico dirigido al equipo médico y
    administrativo durante sus jornadas habituales, asegurando
    el manejo fluido de la herramienta antes de la medición
    final. \\ \hline
Medición final
(post-implemen\-tación)
  & Registro cronometrado del ciclo administrativo completo
    sobre una muestra aleatoria de ciclos de atención, con el
    sistema centralizado en operación real, bajo las mismas
    condiciones y criterios aplicados en la medición de línea
    base. \\ \hline
Comparación cuantitativa
de resultados
  & Contraste estadístico de los tiempos promedio por ciclo
    de atención obtenidos antes y después de la intervención,
    mediante la prueba $t$ de Student para muestras
    relacionadas ($\alpha = 0{,}05$), para determinar si la
    reducción observada es estadísticamente significativa. \\ \hline
\end{tabular}
\caption{Actividades de la Fase 3: Implementación y evaluación
         de impacto.}
\label{tab:actividades_fase3}
\end{table}

\noindent\textbf{Resultado de la Fase 3:} El sistema de
información centralizado operativo en la clínica, acompañado
de la evidencia cuantitativa que comprueba su efectividad en
la reducción del tiempo promedio de gestión administrativa por
ciclo completo de atención.